\documentclass{article}

\usepackage{amsmath, amssymb, amsthm}

\usepackage{logicpuzzle}

% YYYY-MM-DD Date format
\usepackage[style=iso]{datetime2}

% Allows use of internal commands
\makeatletter

% Place an X for an empty cell
\newcommand*\placex[2]{
    \LP@ingrid{#1}{#2}{\LP@NG@columns}{\LP@NG@rows}{nonogram}
    \LP@G@setcellcontent[hcenter,vcenter]{#1}{#2}{$\times$}
    }

% For filling a solved puzzle's empty squares with X marks
% Current solution is to fill every cell with an X and filled cells will cover it.
% Partially finished puzzles will need each desired X placed individually.
% TODO is there anything less stupid
\newcommand{\nonogramfill}[2]{
    \foreach \i in {1,...,#1}
    \foreach \j in {1,...,#2}
    {
        \placex{\i}{\j}
    }
    }


\title{Nonograms}
\author{Michael Gadhia}
\date{\today}

\begin{document}
\maketitle

% Basic explanation and definitions
\section{What Are Nonograms?}

Nonogram puzzles, also called Picross, are puzzles solved by completing a grid of squares.
Using numbers for each row and column, you must determine whether each square is full (which we'll fill in black) or empty (which we'll mark with $\times$).
Often, the end result will be a picture of something!

In each row or column, the numbers (in order) tell you about the filled cells in that row.
For example, a 2 means there's a block of 2 filled cells touching each other in that line.
A 3 and then a 1 means there is a block of 3 filled cells, then some amount of empty cells to separate the blocks, then 1 filled cell on its own.
Let's look at some visual examples.

\begin{center}
    \begin{nonogram}[rows=1,columns=5,scale=0.35,width=3cm,
        fontsize=footnotesize,helplines=5,extracells=2, solution=true]
        \nonogramV{{3,1}}
    \end{nonogram}
\end{center}

There's only one way to fill this row.
The first three cells are the block of 3, then there's room for only one empty space, then the last cell must be the 1.

\begin{center}
    \begin{nonogram}[rows=1,columns=5,scale=0.35,width=3cm,
        fontsize=footnotesize,helplines=5,extracells=2, solution=true]
        \nonogramV{{1,3}}
        \nonogramrow{1}{1/3, 5/1}
        \placex{4}{1}
    \end{nonogram}
\end{center}

Sometimes, we can't immediately solve a line.

\begin{center}
    \begin{nonogram}[rows=1,columns=10,scale=0.35,width=3cm,
        fontsize=footnotesize,helplines=5,extracells=2, solution=true]
        \nonogramV{{5,3}}
    \end{nonogram}
\end{center}

Here, there are three possibilities.
If the block of 5 starts on the first cell, there could be one or two spaces after it, so the block of 3 could start on the seventh or eighth cell.


% Worked example with step-by-step visuals
\section{Solving Nonograms}

% Examples of unique solution, multiple solutions, and no solution
\subsection{Solvability}

% Brief explanation of line-solving techniques
\section{Useful Techniques}



\begin{center}
    % Example puzzle
    \begin{nonogram}[rows=5,columns=5,scale=0.35,width=3cm,
        fontsize=footnotesize,helplines=5,extracells=2]
        % Pay attention to order numbers are defined in, since I don't fully understand it without example

        \nonogramV{{1,1},{4},{1,2},{1},{1}}
        \nonogramH{{1},{4},{1,1},{1},{3}}

    \end{nonogram}
    \hspace{1cm}
    % Solution to adjacent puzzle, but did not set solution=true, need to determine preference for solution visuals 
    \begin{nonogram}[rows=5,columns=5,scale=0.35,width=3cm,
        fontsize=footnotesize,helplines=5,extracells=2]

        \nonogramV{{1,1},{4},{1,2},{1},{1}}
        \nonogramH{{1},{4},{1,1},{1},{3}}

        % nonogramfill command defined above
        \nonogramfill{5}{5}

        \nonogramrow{5}{3/1}
        \nonogramrow{4}{2/1}
        \nonogramrow{3}{1/2, 5/1}
        \nonogramrow{2}{2/4}
        \nonogramrow{1}{2/1, 5/1}
        
    \end{nonogram}
\end{center}


\end{document}